\element{fEtat}{
  \begin{questionmult}{enthalpie}\bareme{haut=2,p=0}
    L'enthalpie
    \begin{multicols}{2}
      \begin{reponses}
        \bonne{se note $H$}
        \bonne{$=U+PV$}
        \mauvaise{A pour variables naturelles $T$ et $P$.}
        \mauvaise{Est constante pour une isobare isotherme}
      \end{reponses}
    \end{multicols}
  \end{questionmult}
}
\element{fEtat}{
  \begin{questionmult}{enthalpieG}\bareme{haut=2,p=0}
    L'enthalpie libre
    \begin{multicols}{2}
      \begin{reponses}
        \bonne{se note $G$}
        \mauvaise{$=U+PV+TS$}
        \bonne{A pour variables naturelles $T$ et $P$.}
        \mauvaise{Est constante pour une isobare adiabatique.}
      \end{reponses}
    \end{multicols}
  \end{questionmult}
}
\setgroupmode{fEtat}{cyclic}

\element{potChim}{
  \begin{questionmult}{potChim}\bareme{haut=2,p=0}
    Le potentiel chimique
    % \begin{multicols}{2}
      \begin{reponses}
        \bonne{dépend de la température et de la pression.}
        \mauvaise{est égal à l'enthalpie molaire.}
        \mauvaise{$=\mu^\circ-RT\ln(a)$}
        \bonne{est le même pour deux phases coexistant à l'équilibre.}
      \end{reponses}
    % \end{multicols}
  \end{questionmult}
}

\element{enthalpieGReaction}{
  \begin{questionmult}{enthalpieGReaction}\bareme{haut=2,p=0}
    L'enthalpie libre de réaction $\Delta_rG$
    % \begin{multicols}{2}
      \begin{reponses}
        \mauvaise{Se mesure en $J$.}
        \mauvaise{$=RT\ln\frac{K^\circ}{Q}$}
        \mauvaise{est positive lorsque la réaction va dans le sens direct.}
        \bonne{est nulle lorsque la réaction est à l'équilibre.}
        \mauvaise{est positive pour une réaction endothermique.}
      \end{reponses}
    % \end{multicols}
  \end{questionmult}
}

\element{vantHoff}{
  \begin{questionmult}{vantHoff}\bareme{haut=2,p=0}
    Selon la relation de Van't Hoff,
    % \begin{multicols}{2}
      \begin{reponses}
        \bonne{$\frac{d\ln K^\circ}{dT}=\frac{\Delta_rH^\circ}{RT^2}$}
        \mauvaise{l'équilibre d'une réaction exothermique se déplace dans le sens direct lorsque la température augmente.}
        \bonne{l'équilibre d'une réaction endothermique se déplace dans le sens inverse lorsque la température diminue.}
        \bonne{il vaut mieux effectuer la synthèse d'une réaction exothermique à basse température.}
        \mauvaise{le système chimique évolue dans le sens qui amplifie les variations de température (loi de Le Chatelier).}
      \end{reponses}
    % \end{multicols}
  \end{questionmult}
}